\documentclass[12pt]{article}
\usepackage[utf8]{inputenc}
\usepackage[spanish]{babel}
\usepackage{geometry}
\usepackage{enumitem}
\usepackage{fancyhdr}

\geometry{a4paper, margin=1in}

\title{Examen de Física Básica: Cinemática y Newton - Forma C}
\author{Nombre: \underline{\hspace{8cm}} \quad Matrícula: \underline{\hspace{3cm}}}
\date{Fecha: \underline{\hspace{4cm}}}

\pagestyle{fancy}
\fancyhf{}
\rhead{Examen Forma C}
\lhead{Física Básica}
\cfoot{\thepage}

\begin{document}

\maketitle

\section*{I. SECCIÓN TEÓRICA (4 Puntos)}

\begin{enumerate}[label=\arabic*.]
    \item \textbf{¿Cuál de las siguientes afirmaciones sobre el desplazamiento y la distancia es correcta?}
    \begin{enumerate}[label=\alph*)]
        \item El desplazamiento siempre es mayor que la distancia.
        \item Son iguales solo si el movimiento es en línea recta y en un solo sentido.
        \item La distancia es una cantidad vectorial.
        \item El desplazamiento es la longitud de la trayectoria.
    \end{enumerate}

    \item \textbf{La Tercera Ley de Newton establece que las fuerzas de acción y reacción:}
    \begin{enumerate}[label=\alph*)]
        \item Actúan sobre el mismo cuerpo y se anulan
        \item Tienen diferentes magnitudes
        \item Actúan sobre cuerpos diferentes y tienen sentidos contrarios
        \item Solo existen en el vacío
    \end{enumerate}

    \item \textbf{La rama de la física que se encarga de la descripción del movimiento, sin considerar las causas que lo producen, es:}
    \begin{enumerate}[label=\alph*)]
        \item Dinámica
        \item Estática
        \item Cinemática
        \item Mecánica Cuántica
    \end{enumerate}

    \item \textbf{En el Movimiento Rectilíneo Uniforme (MRU), ¿cuál de las siguientes características se mantiene constante?}
    \begin{enumerate}[label=\alph*)]
        \item La aceleración
        \item La posición
        \item La velocidad (magnitud, dirección y sentido)
        \item El tiempo
    \end{enumerate}
\end{enumerate}

\section*{II. SECCIÓN PRÁCTICA (6 Puntos)}

\begin{enumerate}[label=\arabic*.]
    \item \textbf{Un tren viaja en línea recta con velocidad constante de 12 m/s durante 5 minutos (300 s). ¿Cuánto se desplazó el tren?}
    \begin{enumerate}[label=\alph*)]
        \item 60 m
        \item 1200 m
        \item 3600 m
        \item 1800 m
    \end{enumerate}

    \item \textbf{Un hombre de 80.0 kg va sentado en un auto que acelera a 2.0 $m/s^2$. La magnitud de la fuerza que el asiento ejerce sobre el hombre es:}
    \begin{enumerate}[label=\alph*)]
        \item 80.0 N
        \item 40.0 N
        \item 160.0 N
        \item 784.0 N
    \end{enumerate}

    \item \textbf{Un pequeño dron se desplaza a una velocidad constante de 36 km/h. ¿Cuál es su velocidad expresada en metros por segundo (m/s)?}
    \begin{enumerate}[label=\alph*)]
        \item 5 m/s
        \item 10 m/s
        \item 15 m/s
        \item 360 m/s
    \end{enumerate}

    \item \textbf{Un objeto parte del reposo y alcanza una velocidad de 50 m/s en 5.0 segundos. ¿Cuál fue su aceleración media?}
    \begin{enumerate}[label=\alph*)]
        \item 5.0 $m/s^2$
        \item 10.0 $m/s^2$
        \item 250 $m/s^2$
        \item 1.0 $m/s^2$
    \end{enumerate}
\end{enumerate}

\end{document}
